% !TeX root = MuJoCo使用指南.tex
% ============================================================
%  第三章　Python 基础用法与可视化
%  内容：MjModel/MjData、仿真推进、状态与控制、命名访问、
%        交互窗口（launch_passive）、离屏渲染、模型转换与保存、常见报错
%  说明：代码清单使用 ASCII 字符，中文说明写在清单之外。
% ============================================================

第二章把环境装好了，本章解决「怎么用」。MuJoCo 的 Python 绑定是\textbf{对 C API 的近乎直接的映射}：所有函数名、参数顺序都与 C 一致，只有少数地方为了符合 Python 习惯做了调整（例如 \texttt{mj\_step} 多了一个 \texttt{nstep} 参数，数组长度参数被省略）。理解这一点，就能直接查官方 C API 文档来写 Python 代码。

\section{模型与数据}

\subsection{两个核心对象}

MuJoCo 把「不变的模型」与「变化的状态」严格分开：

\begin{definition}[mjModel]
	模型的\textbf{编译结果}：几何、关节、质量与惯量、执行器、传感器、选项与默认值。它在仿真过程中\textbf{保持不变}，对应 Python 类 \texttt{mujoco.MjModel}。
\end{definition}

\begin{definition}[mjData]
	\textbf{全部随时间变化的量}与中间计算结果：广义坐标、速度、加速度、控制量、接触力、雅可比、传感器读数等。运行期\textbf{不再分配内存}，全部预分配在 \texttt{MjData} 里。对应 Python 类 \texttt{mujoco.MjData}。
\end{definition}

这种划分带来两个好处：同一个模型可以创建任意多个数据实例（多线程并行采样），而模型本身可以安全共享。

\subsection{加载模型}

\texttt{MjModel} 没有 Python 构造函数，只能通过三个工厂方法创建：

\begin{lstlisting}[style=code, language=Python]
import mujoco

m = mujoco.MjModel.from_xml_path("model/robot.xml")
m = mujoco.MjModel.from_binary_path("model/robot.mjb")
m = mujoco.MjModel.from_xml_string(xml_string, assets_dict)
\end{lstlisting}

\begin{itemize}
	\item \texttt{from\_xml\_path}：从文件加载 MJCF 或 URDF，\textbf{格式由扩展名与内容自动识别}，因此可以直接打开 \texttt{robot.urdf}；
	\item \texttt{from\_binary\_path}：加载编译好的二进制模型（\texttt{.mjb}），加载快但不含 XML 信息、且与版本绑定；
	\item \texttt{from\_xml\_string}：从字符串加载，第二个参数是可选的「虚拟文件系统」字典，把网格、纹理等内容直接以内存里的字节流提供给编译器——\textbf{把模型打包进程序或数据库时非常有用}，也避免了路径问题。
\end{itemize}

\subsection{常用字段}

\begin{table}[H]
	\centering
	\caption{最常用的模型与数据字段}
	\label{tab:Python:常用字段}
	\footnotesize
	\begin{tabular}{>{\raggedright\arraybackslash}p{4.6cm}>{\raggedright\arraybackslash}p{9.6cm}}
		\hline
		字段 & 含义 \\
		\hline
		\texttt{model.nq} & 广义坐标个数（位置） \\
		\texttt{model.nv} & 自由度个数（速度），无自由关节时 $n_{v}=n_{q}$ \\
		\texttt{model.nu} & 执行器个数，即 \texttt{data.ctrl} 的长度 \\
		\texttt{model.nbody} / \texttt{ngeom} / \texttt{njnt} / \texttt{nsite} & 物体、几何体、关节、站点的个数 \\
		\texttt{model.opt.timestep} & 积分步长，默认 $0.002$ s \\
		\texttt{model.body\_mass} & 各物体的质量，模型装错时第一个该查的地方 \\
		\texttt{data.time} & 当前仿真时刻 \\
		\texttt{data.qpos} & 广义坐标，长度 \texttt{nq} \\
		\texttt{data.qvel} & 广义速度，长度 \texttt{nv} \\
		\texttt{data.qacc} & 广义加速度，长度 \texttt{nv} \\
		\texttt{data.ctrl} & 控制输入，长度 \texttt{nu} \\
		\texttt{data.act} & 执行器内部状态（如 \texttt{pid}、\texttt{dcmotor} 的内置状态） \\
		\texttt{data.sensordata} & 传感器读数，长度由模型决定 \\
		\texttt{data.xpos} / \texttt{data.xquat} & 各物体在世界坐标系中的位置与姿态（$n_{q}$ 的位形经正运动学算出） \\
		\texttt{data.qfrc\_applied} / \texttt{data.xfrc\_applied} & 用户施加的广义力 / 笛卡尔力，用来做外部扰动 \\
		\hline
	\end{tabular}
\end{table}

\begin{warnbox}[数组是「视图」而不是「副本」]
	\texttt{data.qpos} 返回的是指向引擎内存的 NumPy 视图。写 \texttt{history.append(data.qpos)} 会得到一个「所有元素都相同」的列表——因为每一步 \texttt{mj\_step} 都会就地改写这块内存。记录数据时必须 \texttt{.copy()}：
	\texttt{history.append(data.qpos.copy())}。
\end{warnbox}

\section{推进仿真}

\subsection{三个入口}

\begin{lstlisting}[style=code, language=Python]
mujoco.mj_step(m, d)              # 前进一个时间步（完整的前向动力学 + 积分）
mujoco.mj_step(m, d, nstep=20)    # Python 专有：连续走 20 步，中途不回到解释器
mujoco.mj_forward(m, d)           # 只做前向计算，不推进时间（改完状态后先调它）
mujoco.mj_step1(m, d); ...        # 只算到「位置相关」的部分
mujoco.mj_step2(m, d)             # 只算到「速度相关」的部分
mujoco.mj_resetData(m, d)         # 回到模型的初始状态
\end{lstlisting}

\texttt{mj\_step1}/\texttt{mj\_step2} 把一步拆成两半，中间是\textbf{施加控制的正确位置}（\texttt{mjcb\_control} 回调即在此处被调用）。若只是「推进仿真」，用 \texttt{mj\_step} 即可；想自己写控制器且不装回调，用下面的顺序更清晰：

\begin{lstlisting}[style=code, language=Python]
mujoco.mj_step1(m, d)
d.ctrl[:] = policy(d)       # 或直接改 d.qfrc_applied
mujoco.mj_step2(m, d)
\end{lstlisting}

\subsection{状态存取与关键帧}

\begin{lstlisting}[style=code, language=Python]
spec = mujoco.mjtState.mjSTATE_FULLPHYSICS
n = mujoco.mj_stateSize(m, spec)
buf = np.empty(n)
mujoco.mj_getState(m, d, buf, spec)     # 存
mujoco.mj_setState(m, d, buf, spec)     # 取
mujoco.mj_resetDataKeyframe(m, d, 0)    # 回到第 0 个关键帧 <key .../>
\end{lstlisting}

整状态的存取比逐个变量抄写更安全，做强化学习环境的重置（reset）时尤其有用。若模型里定义了 \texttt{<key>}（关键帧），直接用 \texttt{mj\_resetDataKeyframe} 即可回到一个给定的初始姿态。

\subsection{一个完整的受控例子}

下面给方块加一个滑轨与一个位置执行器，并用简单比例控制把它拉到目标高度。把它和第三章开头的 \texttt{hello.xml} 对照着看，可以体会 MJCF 与 Python 两侧的分工：\textbf{模型里定义「有什么」，程序里定义「怎么用」}。

\begin{lstlisting}[style=xml]
<mujoco>
  <option timestep="0.002"/>
  <worldbody>
    <light pos="0 0 3" dir="0 0 -1"/>
    <geom type="plane" size="1 1 0.1"/>
    <body name="slider" pos="0 0 0.5">
      <joint name="slide" type="slide" axis="0 0 1"/>
      <geom type="box" size=".05 .05 .05" rgba="0 .6 .9 1"/>
    </body>
  </worldbody>
  <actuator>
    <position name="lift" joint="slide" kp="200" kv="20"/>
  </actuator>
</mujoco>
\end{lstlisting}

\begin{lstlisting}[style=code, language=Python]
import mujoco

m = mujoco.MjModel.from_xml_path("slider.xml")
d = mujoco.MjData(m)

target = 0.8
act_id = m.actuator("lift").id

for _ in range(2000):
    d.ctrl[act_id] = target
    mujoco.mj_step(m, d)

print("time =", d.time)
print("height =", d.joint("slide").qpos[0])
\end{lstlisting}

期望输出中 \texttt{height} 收敛到 $0.8$ 附近。注意这里用的是 MJCF 自带的位置执行器（内部就是 PD），比自己写 PD 更省事，也更容易复现。

\section{命名访问}

\subsection{按名字拿数组}

MuJoCo 的 C API 用整数 id 索引数组，Python 绑定在此之上提供了「命名访问」：\texttt{m.类别('名字')} 返回一个访问器，属性名就是去掉前缀的字段名，并且\textbf{访问是 $O(1)$ 的，不随模型规模变慢}。

\begin{lstlisting}[style=code, language=Python]
m.geom("gizmo").rgba            # 长度为 4 的 NumPy 视图，可直接赋值改颜色
m.geom("gizmo").id              # 等价于 mujoco.mj_name2id(m, mjOBJ_GEOM, "gizmo")
m.geom(3).name                  # 反向查询
d.body("slider").xpos           # 物体的世界坐标
d.joint("slide").qpos           # 只属于该关节的那一段 qpos
m.joint("slide").range          # 关节限位
\end{lstlisting}

可用的类别名见表~\ref{tab:Python:命名类别}。关节类别同时接受 \texttt{jnt} 与 \texttt{joint}，相机接受 \texttt{cam} 与 \texttt{camera}，其余别名同理。

\begin{table}[H]
	\centering
	\caption{命名访问的类别}
	\label{tab:Python:命名类别}
	\footnotesize
	\begin{tabular}{lllll}
		\hline
		\texttt{body} & \texttt{jnt/joint} & \texttt{geom} & \texttt{site} & \texttt{cam/camera} \\
		\texttt{light} & \texttt{mesh} & \texttt{skin} & \texttt{hfield} & \texttt{tex/texture} \\
		\texttt{mat/material} & \texttt{pair} & \texttt{exclude} & \texttt{eq/equality} & \texttt{tendon/ten} \\
		\texttt{actuator} & \texttt{sensor} & \texttt{numeric} & \texttt{text} & \texttt{key/keyframe} \\
		\hline
	\end{tabular}
\end{table}

\begin{tipbox}[工程习惯]
	凡是程序里要反复引用的元素——关节、执行器、站点、相机——\textbf{都在 MJCF 里起名字}，然后在 Python 里用命名访问取用。这样模型改名或增删部件时，程序里对应的代码不必跟着改 id。
\end{tipbox}

\section{交互式可视化}

\subsection{三种用法}

官方在 \texttt{mujoco.viewer} 模块里提供三种用法（表~\ref{tab:Python:viewer}）。它们的区别只在「谁负责推进物理」：managed 模式由查看器自己推进，passive 模式由你的脚本推进。

\begin{table}[H]
	\centering
	\caption{三种查看器用法}
	\label{tab:Python:viewer}
	\footnotesize
	\begin{tabular}{>{\raggedright\arraybackslash}p{2.6cm}>{\raggedright\arraybackslash}p{5.0cm}>{\raggedright\arraybackslash}p{6.6cm}}
		\hline
		用法 & 调用 & 适用 \\
		\hline
		managed & \texttt{viewer.launch(model, data)} & 只看模型、手动拖动交互；\textbf{会阻塞你的代码} \\
		standalone & \texttt{python -m mujoco.viewer --mjcf=...} & 命令行直接看模型文件，等价于 C++ 的 \texttt{simulate} \\
		passive & \texttt{viewer.launch\_passive(model, data)} & 自己写控制循环，同时把画面显示出来；\textbf{不阻塞} \\
		\hline
	\end{tabular}
\end{table}

\subsection{passive 查看器}

这是做控制与学习最常用的模式。要点有三条：\textbf{用 \texttt{with} 管理生命周期}、\textbf{每步调用 \texttt{sync()}}、\textbf{改动物理或可视化状态前先 \texttt{lock()}}。

\begin{lstlisting}[style=code, language=Python]
import time
import mujoco
import mujoco.viewer

m = mujoco.MjModel.from_xml_path("slider.xml")
d = mujoco.MjData(m)

paused = False

def key_callback(keycode):
    global paused
    if chr(keycode) == " ":
        paused = not paused

with mujoco.viewer.launch_passive(m, d, key_callback=key_callback) as viewer:
    start = time.time()
    while viewer.is_running() and time.time() - start < 30:
        step_start = time.time()

        if not paused:
            d.ctrl[0] = 0.8
            mujoco.mj_step(m, d)

            with viewer.lock():
                viewer.opt.flags[mujoco.mjtVisFlag.mjVIS_CONTACTPOINT] = int(d.time % 2)

            viewer.sync()

        dt = m.opt.timestep - (time.time() - step_start)
        if dt > 0:
            time.sleep(dt)
\end{lstlisting}

\begin{itemize}
	\item \texttt{sync()} 是「把 \texttt{mjModel}/\texttt{mjData} 交给 GUI，再把 GUI 上的改动（鼠标扰动、控制输入、各种开关）取回来」的唯一时机。\textbf{不调用 \texttt{sync} 窗口里的东西不会动}。
	\item \texttt{sync(state\_only=True)} 只同步状态相关的字段并做一次 \texttt{mj\_forward}，快得多；但它不会反映你对 \texttt{mjModel} 的改动（比如改颜色）。
	\item \texttt{viewer.cam}、\texttt{viewer.opt}、\texttt{viewer.pert} 分别对应相机、可视化选项与扰动对象，可以直接改。
	\item 想要「按真实时间跑」，用上面末尾的 \texttt{time.sleep} 补足每步耗时；纯追求算得快就把它去掉。
	\item \texttt{key\_callback} 收到的是键码，\texttt{chr()} 可转成字符；空格暂停是最常用的例子。
\end{itemize}

\begin{warnbox}[平台注意]
	Windows 与 Linux 上 \texttt{launch\_passive} 直接在普通 \texttt{python} 下运行即可，但\textbf{必须在主线程}调用。
	macOS 上则必须用包内附带的 \texttt{mjpython} 启动脚本（\texttt{mjpython my\_script.py}），否则会因为系统要求「渲染必须在主线程」而失败。
\end{warnbox}

\subsection{往画面里加自定义图形}

\texttt{viewer.user\_scn} 是一个完全由你控制的 \texttt{mjvScene}，可以在里面追加几何体，用来画轨迹、力箭头、目标点等：

\begin{lstlisting}[style=code, language=Python]
import numpy as np
import mujoco

with mujoco.viewer.launch_passive(m, d) as viewer:
    while viewer.is_running():
        mujoco.mj_step(m, d)

        viewer.user_scn.ngeom = 0
        mujoco.mjv_initGeom(
            viewer.user_scn.geoms[0],
            type=mujoco.mjtGeom.mjGEOM_SPHERE,
            size=[0.02, 0, 0],
            pos=0.1 * np.array([1, 0, 2]),
            mat=np.eye(3).flatten(),
            rgba=np.array([1.0, 0.2, 0.2, 1.0]),
        )
        viewer.user_scn.ngeom = 1

        viewer.sync()
\end{lstlisting}

\section{离屏渲染与录像}

做论文插图、视觉伺服、或把仿真接到图像模型上时，需要\textbf{不开窗口}地渲染出图像。

\subsection{基本用法}

\texttt{mujoco.Renderer} 是对底层 OpenGL 上下文（\texttt{mujoco.GLContext}）的封装：创建时指定分辨率，之后每帧先 \texttt{update\_scene} 再 \texttt{render}。

\begin{lstlisting}[style=code, language=Python]
import mujoco
import numpy as np
from PIL import Image

m = mujoco.MjModel.from_xml_path("slider.xml")
d = mujoco.MjData(m)

renderer = mujoco.Renderer(m, height=480, width=640)
try:
    for _ in range(500):
        mujoco.mj_step(m, d)

    renderer.update_scene(d)
    rgb = renderer.render()
    print(rgb.shape, rgb.dtype)
    Image.fromarray(rgb).save("frame.png")

    renderer.enable_depth_rendering()
    renderer.update_scene(d)
    depth = renderer.render()
    print(depth.shape, depth.dtype, depth.min(), depth.max())
    renderer.disable_depth_rendering()
finally:
    renderer.close()
\end{lstlisting}

\begin{itemize}
	\item \texttt{update\_scene} 可以用 \texttt{camera=} 指定模型里定义的相机（\texttt{<camera name="cam"/>}），也可以传 \texttt{scene\_option} 控制显示哪些元素：
	      \texttt{renderer.update\_scene(d, camera="cam", scene\_option=opt)}。
	\item 深度渲染给出的深度值不是「米」，其换算与相机的 \texttt{fovy}、\texttt{znear}、\texttt{zfar} 有关，做视觉任务时建议自己按近远平面做一次线性反解，不要直接当距离用。
	\item 渲染必须在\textbf{同一个线程}里进行：OpenGL 上下文一次只能在一个线程上生效。
\end{itemize}

\subsection{Windows 上的后端限制}

MuJoCo 的离屏渲染在不同平台依赖不同的 OpenGL 上下文实现：Linux 支持 GLX、OSMesa 与 EGL，而 \textbf{Windows 与 macOS 使用操作系统自带的 OpenGL 实现}，也就是必须有一个（可以是隐藏的）窗口上下文。因此：

\begin{itemize}
	\item Windows 上\textbf{不要}照搬 Linux 教程里的 \texttt{MUJOCO\_GL=egl} 或 \texttt{MUJOCO\_GL=osmesa}，这两个后端在 Windows 上不可用；
	\item 无显卡的服务器上跑 Windows 虚拟机时，离屏渲染会因为只有 Microsoft 基础显示驱动（OpenGL 1.1）而失败；此时应改用 Linux 容器 + EGL，或者退回到纯计算（不渲染）。
\end{itemize}

\begin{tipbox}[录像]
	把每帧的 \texttt{rgb} 收集起来写入视频即可，例如用 \texttt{imageio}：
	\texttt{imageio.mimsave("run.mp4", frames, fps=60)}。注意每帧都要先推进物理，再 \texttt{update\_scene} 后 \texttt{render}。
\end{tipbox}

\section{模型转换与保存}

建模阶段常常需要「加载一个 URDF、稍作修改、再存成 MJCF」。旧接口是「保存最近一次编译的 XML」：

\begin{lstlisting}[style=code, language=Python]
import mujoco

m = mujoco.MjModel.from_xml_path("robot.urdf")   # 直接编译 URDF
mujoco.mj_saveLastXML("robot_from_urdf.xml")     # 保存最近一次编译的 XML
\end{lstlisting}

新接口围绕 \texttt{MjSpec}（与 MJCF 一一对应的可编辑对象）展开，适合在程序里改模型，也适合连同资源一起打包：

\begin{lstlisting}[style=code, language=Python]
import mujoco

spec = mujoco.MjSpec.from_file("robot.urdf")
spec.modelname = "two_link_arm"
model = spec.compile()
spec.to_zip("robot.mjz")               # XML + 网格 + 纹理 打包成一个文件
spec2 = mujoco.MjSpec.from_zip("robot.mjz")
\end{lstlisting}

\texttt{.mjz} 是一个 zip 归档，模型、网格与纹理全在里面，给别人或换机器时不会出现「找不到 \texttt{mesh} 文件」的问题，适合做交付。第四章的 C++ 侧同样可以直接加载 \texttt{.mjz}。

\section{常见运行期问题}

\begin{table}[H]
	\centering
	\caption{Python 运行期常见问题}
	\label{tab:Python:运行报错}
	\footnotesize
	\begin{tabular}{>{\raggedright\arraybackslash}p{4.6cm}>{\raggedright\arraybackslash}p{4.4cm}>{\raggedright\arraybackslash}p{5.4cm}}
		\hline
		现象 & 原因 & 处理 \\
		\hline
		\texttt{ValueError: Error: could not open file ...} & 路径不对，或 XML 里的 \texttt{mesh}/\texttt{texture} 找不到 & 用绝对路径试一次；确认 \texttt{meshdir} 与文件实际位置一致 \\
		\texttt{ValueError: Error: unknown element ...} & MJCF 拼写错误 & MuJoCo 的解析器报错信息通常给出出错的行号与元素名，直接照行号看 \\
		\texttt{mujoco.FatalError} & 引擎的不可恢复错误（如数值发散后的断言） & 先降到更小的时间步长；检查质量与惯量是否合理 \\
		模型加载后什么也看不见 & XML 里没有 \texttt{<light>}，或缺相机 & 用 \texttt{simulate} 或 \texttt{viewer} 打开同一文件对比；补光源 \\
		窗口一闪而过或打不开 & 远程桌面、显卡驱动过旧 & 参见 1.3 节的说明 \\
		打印的状态全都一样 & 记录了 NumPy 视图而非副本 & 记住 \texttt{.copy()} \\
		仿真「飞掉」 & 时间步过大、惯量不合理、初始穿插 & 见第九章的调参与排错 \\
		\hline
	\end{tabular}
\end{table}
